% =============================================================================
%  Section 4.2 — Thực nghiệm
%  Subsections:
%    4.2.1  Mục tiêu và câu hỏi nghiên cứu
%    4.2.2  Kiểm chứng tính hợp lệ của CheoBench v2
%    4.2.3  Độ đo đánh giá
%    4.2.4  Kết quả thực nghiệm
%    4.2.5  Phân tích và thảo luận
% =============================================================================

\section{Thực nghiệm}
\label{sec:experiments}

\subsection{Mục tiêu và câu hỏi nghiên cứu}

Thực nghiệm nhằm đánh giá một cách hệ thống hiệu quả của kiến trúc GraphRAG đề xuất đối với việc khai thác tri thức có cấu trúc trong miền nghệ thuật Chèo, đồng thời phân tích đóng góp của từng thành phần trong kiến trúc. Bốn câu hỏi nghiên cứu được xác lập:

\begin{enumerate}
  \item Mức độ truy xuất trong Knowledge Graph ảnh hưởng như thế nào đến hiệu năng tổng thể?
  \item Định dạng biểu diễn ngữ cảnh tác động ra sao đến khả năng suy luận của mô hình ngôn ngữ?
  \item Chiến lược tạo sinh có cải thiện tính trung thực và độ chính xác của câu trả lời hay không?
  \item So với các hệ thông RAG truyền thống và LLM-only, GraphRAG mang lại lợi ích gì trong miền tri thức có cấu trúc quan hệ phức tạp?
\end{enumerate}

% Cần chỉnh sửa thêm
\subsection{Kiểm chứng tính hợp lệ của CheoBench v2}
\label{sec:benchmark_validity}

Tính hợp lệ của bộ dữ liệu được đánh giá trên bốn chiều theo khung Messick~\cite{messick1995validity}: bao phủ nội dung, cấu trúc độ khó, khả năng phân biệt hệ thống và nhất quán của hệ độ đo. Số liệu được tính từ bộ câu hỏi và từ kết quả 100 câu ở Mục~\ref{sec:experiments}.

Về bao phủ nội dung, 100 câu hỏi liệt kê 225 thực thể trong trường \\ \texttt{related\_entities}; đối chiếu với Knowledge Graph, bộ câu hỏi phủ 7/7 vở chèo, 25/34 nhân vật, 23/41 diễn viên và 8/15 trích đoạn (Bảng~\ref{tab:coverage_entity}). Toàn bộ 225 thực thể đều khớp với định danh trong KG, không câu nào yêu cầu tri thức ngoài phạm vi hệ thống. Coverage ở lớp vở đạt 100\%; các lớp con dao động từ 53\% đến 73\%, tập trung vào các thực thể trung tâm theo thực hành của BEIR~\cite{thakur2021beir}. 100 câu hỏi được phân vào 10 subcategory khác nhau, trải trên nhiều dạng truy vấn.

\begin{table}[h]
\centering
\caption{Mức độ bao phủ của CheoBench v2 trên Knowledge Graph Chèo.}
\label{tab:coverage_entity}
\begin{tabular}{lrrr}
\hline
\textbf{Loại thực thể} & \textbf{Đã phủ} & \textbf{Tổng} & \textbf{Tỷ lệ} \\
\hline
Vở chèo       &  7 &  7 & 100.0\% \\
Nhân vật      & 25 & 34 &  73.5\% \\
Diễn viên     & 23 & 41 &  56.1\% \\
Trích đoạn    &  8 & 15 &  53.3\% \\
\hline
\end{tabular}
\end{table}

Về cấu trúc độ khó, Edge et al.~\cite{edge2024graphrag} dự đoán khoảng cách GraphRAG--RAG giảm dần từ truy vấn Local (ánh xạ tên riêng) đến truy vấn Global (tổng hợp toàn cục). Bảng~\ref{tab:construct_validity} xác nhận xu hướng này trên 100 câu hỏi: khoảng cách giảm đơn điệu từ +0,309 (Local) qua +0,245 (Community) xuống +0,163 (Global). Số thực thể liên quan trung bình cũng phân tách ba nhóm ở mức 1,94 (Local), 2,50 (Community) và 2,33 (Global).

\begin{table}[h]
\centering
\caption{Khoảng cách hiệu năng GraphRAG--RAG theo nhóm truy vấn.}
\label{tab:construct_validity}
\begin{tabular}{lcccc}
\hline
\textbf{Nhóm} & \textbf{\#Ents TB} & \textbf{GraphRAG} & \textbf{RAG} & \textbf{Khoảng cách} \\
\hline
Local      & 1.94 & 0.862 & 0.553 & +0.309 \\
Community  & 2.50 & 0.861 & 0.616 & +0.245 \\
Global     & 2.33 & 0.782 & 0.619 & +0.163 \\
\hline
\end{tabular}
\end{table}

Về khả năng phân biệt hệ thống, trên 100 cặp (GraphRAG, RAG) không có cặp hoà điểm, 94 cặp GraphRAG đạt điểm cao hơn, $\Delta = S_{\text{GraphRAG}} - S_{\text{RAG}}$ có trung bình +0,240 và độ lệch chuẩn 0,193. Phân phối trải rộng với 72\% cặp có $|\Delta| \geq 0{,}10$ và 33\% cặp có $|\Delta| \geq 0{,}30$ (Bảng~\ref{tab:discrimination}), cho phép phân biệt hai hệ thống ở nhiều cấp độ thay vì ở hai cực.

\begin{table}[h]
\centering
\caption{Phân phối khoảng cách điểm tổng hợp $\Delta = S_{\text{GraphRAG}} - S_{\text{RAG}}$ trên 100 câu hỏi.}
\label{tab:discrimination}
\begin{tabular}{lrr}
\hline
\textbf{Chỉ số} & \textbf{Giá trị} & \textbf{Tỷ lệ} \\
\hline
Câu có GraphRAG $\geq$ RAG & 94/100 & 94.0\% \\
Câu hoà điểm ($|\Delta| < 10^{-6}$) &  0/100 &  0.0\% \\
Câu có $|\Delta| \geq 0.05$ & 84/100 & 84.0\% \\
Câu có $|\Delta| \geq 0.10$ & 72/100 & 72.0\% \\
Câu có $|\Delta| \geq 0.20$ & 50/100 & 50.0\% \\
Câu có $|\Delta| \geq 0.30$ & 33/100 & 33.0\% \\
\hline
Trung bình $\Delta$ & +0.240 & --- \\
Độ lệch chuẩn của $\Delta$ &  0.193 & --- \\
\hline
\end{tabular}
\end{table}

Về nhất quán của hệ độ đo, Cronbach's $\alpha$~\cite{cronbach1951alpha} tính trên điểm per-case của GraphRAG cho khối IR năm độ đo (MAP, NDCG@10, Precision, Recall, ContextEntitiesRecall) đạt 0,790, vượt ngưỡng 0,70~\cite{nunnally1978psychometric}; $\alpha$ toàn cục chín độ đo là 0,699 (Bảng~\ref{tab:consistency}). Hai khối Context và Generation có $\alpha$ gần 0, cho thấy các cặp độ đo trong mỗi khối đo các khía cạnh khác nhau (độ liên quan so với độ bao phủ; độ trung thực so với độ phù hợp). Tương quan Pearson bổ sung: $r(\text{MAP}, \text{NDCG@10}) = 0{,}98$, Faithfulness trực giao với các độ đo khác ($|r| \leq 0{,}06$), ContextEntitiesRecall tương quan yếu với mọi độ đo ($|r| \leq 0{,}14$). Bộ chín độ đo do đó vừa có khối nhất quán nội tại, vừa có các chiều bổ sung trực giao.

\begin{table}[h]
\centering
\caption{Cronbach's $\alpha$ của các khối độ đo trên 100 câu hỏi (hệ thống GraphRAG).}
\label{tab:consistency}
\begin{tabular}{lcc}
\hline
\textbf{Khối độ đo} & \textbf{Số độ đo} & \textbf{Cronbach's $\alpha$} \\
\hline
Toàn bộ                   & 9 & 0.699 \\
Khối IR xác định          & 5 & \textbf{0.790} \\
Khối Context (RAGAs)      & 2 & $-0.048$ \\
Khối Generation (RAGAs)   & 2 & $-0.063$ \\
\hline
\end{tabular}
\end{table}


\subsection{Độ đo đánh giá}

\subsubsection{Chất lượng truy xuất}

Chất lượng truy xuất được đánh giá bởi bốn độ đo IR truyền thống kết hợp với ba độ đo ngữ nghĩa theo RAGAs~\cite{es2023ragas}. Nhóm IR gồm:

\begin{itemize}
  \item \textbf{MAP}~\cite{manning2008ir} --- trung bình độ chính xác (average precision) trên toàn bộ truy vấn; trong bài toán này, MAP thưởng các cấu hình đưa nhiều thực thể đúng lên vị trí đầu của danh sách truy xuất, phản ánh chất lượng xếp hạng tổng thể.

  \item \textbf{NDCG@10}~\cite{jarvelin2002ndcg} --- đánh giá chất lượng xếp hạng ở mười vị trí đầu với phạt logarit theo vị trí; mức cắt $K = 10$ được chọn để phù hợp với kích thước ngữ cảnh điển hình mà pipeline đưa vào mô hình ngôn ngữ.

  \item \textbf{Precision} và \textbf{Recall} --- tỷ lệ chính xác và độ bao phủ của kết quả truy xuất so với tập thực thể liên quan $R$ theo ground truth, tính trên toàn bộ danh sách truy xuất $L$ sau khi loại trùng lặp:
    \begin{equation}
      \text{Precision} = \frac{|L \cap R|}{|L|}, \qquad \text{Recall} = \frac{|L \cap R|}{|R|}.
    \end{equation}
\end{itemize}

Ba độ đo ngữ nghĩa theo RAGAs \cite{es2023ragas} bổ sung cho các độ đo IR ở khía cạnh ngữ nghĩa và đều nằm trong thang $[0, 1]$. \textbf{Context Precision} và \textbf{Context Recall} được mô hình ngôn ngữ chấm điểm từ cặp (câu hỏi, ngữ cảnh) và (ground truth, ngữ cảnh) để ước lượng tỷ lệ thông tin liên quan và mức độ bao phủ tri thức cần thiết. \textbf{Context Entities Recall} được tính trực tiếp: gọi $E$ là tập các thực thể bắt buộc theo ground truth, độ đo là tỷ lệ thực thể trong $E$ xuất hiện trong ngữ cảnh được truy xuất:

\begin{equation}
\text{Context Entities Recall} = \frac{|E \cap \text{context}|}{|E|}.
\end{equation}

Điểm truy xuất tổng hợp được tính theo công thức:

\begin{equation}
\begin{aligned}
S_{\text{retrieval}} =\;& 0.20 \cdot \text{MAP}
  + 0.15 \cdot \text{NDCG@10} \\
& + 0.10 \cdot \text{Precision}
  + 0.10 \cdot \text{Recall} \\
& + 0.20 \cdot \text{Context Precision}
  + 0.15 \cdot \text{Context Recall} \\
& + 0.10 \cdot \text{Context Entities Recall}.
\end{aligned}
\end{equation}

Trước hết, khối IR được đánh trọng số 0,55 - cao hơn khối RAGAs (0,45), do các độ đo IR có thể tính toán trực tiếp và tái lập chính xác, trong khi RAGAs phụ thuộc vào phán đoán của mô hình ngôn ngữ (LLM-judge) và có sự biến động giữa các lần chạy. Trong khối IR, MAP giữ trọng số cao nhất (0,20) do bao quát toàn bộ chất lượng xếp hạng trên mọi vị trí, NDCG@10 đứng thứ hai (0,15) do chỉ nhạy với mười vị trí đầu, còn Precision và Recall mỗi độ đo chỉ phản ánh một khía cạnh riêng - độ chính xác và độ bao phủ - nên được cân bằng ở 0,10.Trong khối RAGAs, Context Precision được ưu tiên cao nhất (0,20) vì ngữ cảnh chứa thông tin không liên quan dễ dẫn mô hình ngôn ngữ sai hướng và ảnh hưởng trực tiếp đến chất lượng câu trả lời, trong khi Context Recall (0,15) và Context Entities Recall (0,10) cùng đo độ bao phủ nhưng ở các phạm vi khác nhau, với Entities Recall hẹp hơn nên nhận trọng số thấp hơn.

Cơ chế tổng hợp xử lý trường hợp độ đo không khả dụng bằng tái chuẩn hoá. Khi một độ đo RAGAs trả về giá trị vô hiệu do lỗi tạm thời của mô hình, độ đo đó được loại khỏi công thức và các trọng số còn lại được chuẩn hoá lại về tổng 1. Điểm truy xuất do đó luôn được tính trên tập độ đo khả dụng cho từng câu hỏi, tránh việc lỗi kỹ thuật đơn lẻ làm giảm kết quả đánh giá.

\subsubsection{Chất lượng tạo sinh}

Chất lượng tạo sinh được đánh giá qua hai độ đo RAGAs~\cite{es2023ragas}, cả hai đều thuộc loại \textbf{LLM-as-judge} - mô hình ngôn ngữ đóng vai trò bên chấm điểm và trả về giá trị trong thang $[0, 1]$. \textbf{Faithfulness} đo tỷ lệ nội dung câu trả lời thực sự được ngữ cảnh truy xuất hỗ trợ, được thiết kế để phát hiện hallucination --- hiện tượng mô hình ngôn ngữ đưa ra thông tin không có trong ngữ cảnh. Giá trị của độ đo được tính qua hai bước, mỗi bước là một lần gọi mô hình ngôn ngữ:

\begin{enumerate}
  \item \textit{Tách mệnh đề.} Mô hình được yêu cầu liệt kê các khẳng định đơn lẻ trong câu trả lời, mỗi khẳng định trên một dòng. Gọi tập các mệnh đề thu được là $C = \{c_1, \ldots, c_n\}$. Ví dụ, câu trả lời ``Thị Kính là vợ của Thiện Sỹ và bị Sùng Bà vu oan'' được tách thành hai mệnh đề độc lập ``Thị Kính là vợ của Thiện Sỹ'' và ``Thị Kính bị Sùng Bà vu oan''.

  \item \textit{Kiểm chứng từng mệnh đề.} Với mỗi $c_i \in C$, mô hình được đưa ngữ cảnh truy xuất cùng với mệnh đề và được hỏi mệnh đề đó có được hỗ trợ hay không. Mệnh đề chỉ được tính là ``được hỗ trợ'' khi mô hình trả lời ``có''; các trường hợp ``không'' hoặc mơ hồ đều được coi là không hỗ trợ.
\end{enumerate}

Giá trị Faithfulness cuối cùng là tỷ lệ mệnh đề được hỗ trợ trên tổng số mệnh đề trích xuất được:

\begin{equation}
\text{Faithfulness} = \frac{|\{c \in C : \text{được ngữ cảnh hỗ trợ}\}|}{|C|}.
\end{equation}

\textbf{Answer Relevance} được mô hình ngôn ngữ ước lượng trực tiếp từ cặp (câu hỏi, câu trả lời), cho một điểm số liên tục phản ánh mức độ câu trả lời giải quyết đúng câu hỏi được đặt ra. Điểm tạo sinh tổng hợp hai độ đo với trọng số bằng nhau:

\begin{equation}
S_{\text{generation}} = 0.5 \cdot \text{Faithfulness} + 0.5 \cdot \text{Answer Relevance}.
\end{equation}

\subsubsection{Điểm tổng hợp}

Điểm tổng hợp $S_{\text{overall}}$ kết hợp $S_{\text{retrieval}}$ và $S_{\text{generation}}$ với trọng số bằng nhau, đóng vai trò chỉ số chính trong phần kết quả thực nghiệm (Mục~\ref{sec:experiments}). Công thức được định nghĩa như sau:

\begin{equation}
S_{\text{overall}} = 0.5 \cdot S_{\text{retrieval}} + 0.5 \cdot S_{\text{generation}}.
\end{equation}

Cơ chế tái chuẩn hoá ở cấp này hoạt động tương tự các khối con: khi $S_{\text{retrieval}}$ hoặc $S_{\text{generation}}$ không tính được, điểm còn lại được dùng thay thế, đảm bảo mọi câu hỏi trong bộ đánh giá đều có điểm tổng hợp hợp lệ nếu ít nhất một trong hai thành phần thành công.

\subsection{Kết quả thực nghiệm}

\subsubsection{Ảnh hưởng của mức độ truy xuất}

Bảng~\ref{tab:granularity_results} cho thấy hiệu năng tăng đơn điệu khi chuyển từ các mức độ truy xuất đơn giản sang phức tạp hơn. Truy xuất ở mức nút đạt điểm tổng hợp thấp nhất (0,54) do thông tin thuộc tính đơn lẻ không đủ cho các câu hỏi có yếu tố quan hệ. Mức triplet (0,63) và path (0,61) cải thiện nhờ bắt đầu khai thác quan hệ giữa các thực thể, tuy phạm vi ngữ cảnh vẫn còn giới hạn. Mức đồ thị con tăng rõ rệt lên (0,72) khi cung cấp vùng lân cận đủ rộng cho suy luận, và cấu hình kết hợp đa mức độ đạt cao nhất (0,76) - xác nhận các mức độ có tính bổ sung thay vì thay thế lẫn nhau. Chênh lệch giữa kết hợp và đồ thị con đơn lẻ chỉ 0,04 điểm, nhỏ hơn đáng kể các bước nhảy 0,09--0,18 ở những mức thấp, phản ánh hiện tượng lợi ích biên giảm dần khi tăng độ phong phú của ngữ cảnh.

\begin{table}[h]
\centering
\caption{Ảnh hưởng của mức độ truy xuất tới hiệu năng hệ thống.}
\label{tab:granularity_results}
\begin{tabular}{lccc}
\hline
\textbf{Mức độ truy xuất} & \textbf{Retrieval} & \textbf{Generation} & \textbf{Overall} \\
\hline
Node              & 0.46 & 0.62 & 0.54 \\
Triplet           & 0.56 & 0.70 & 0.63 \\
Path              & 0.54 & 0.68 & 0.61 \\
Subgraph          & 0.68 & 0.75 & 0.72 \\
Kết hợp đa mức độ & 0.73 & 0.78 & 0.76 \\
\hline
\end{tabular}
\end{table}

\subsubsection{Ảnh hưởng của định dạng biểu diễn ngữ cảnh}

Bảng~\ref{tab:format_results} cho thấy JSON có cấu trúc đạt điểm tổng hợp cao nhất trong các định dạng đơn lẻ với 0,72, sát ngay trên văn bản tự nhiên 0,71; trong khi bảng kề chỉ đạt 0,68 và chuỗi nút thấp nhất ở 0,65. Kết quả này khẳng định mô hình ngôn ngữ hoạt động tốt hơn khi tri thức được trình bày ở dạng có cấu trúc rõ ràng hoặc gần biểu diễn ngôn ngữ tự nhiên; chuỗi nút tuyến tính do thiếu thông tin quan hệ đầy đủ làm giảm rõ khả năng suy luận. Kết hợp nhiều định dạng đạt 0,76, cao hơn định dạng tốt nhất đơn lẻ 0,04 điểm --- một mức cải thiện khiêm tốn phản ánh cùng hiện tượng lợi ích biên giảm dần khi tăng kích thước ngữ cảnh.

\begin{table}[h]
\centering
\caption{Ảnh hưởng của định dạng biểu diễn ngữ cảnh tới hiệu năng hệ thống.}
\label{tab:format_results}
\begin{tabular}{lc}
\hline
\textbf{Định dạng} & \textbf{Điểm tổng hợp} \\
\hline
Văn bản tự nhiên        & 0.71 \\
Cấu trúc JSON           & 0.72 \\
Bảng kề                 & 0.68 \\
Chuỗi nút               & 0.65 \\
Kết hợp nhiều định dạng & 0.76 \\
\hline
\end{tabular}
\end{table}

\subsubsection{Ảnh hưởng của chiến lược tạo sinh}

So sánh ba chiến lược tạo sinh cho thấy Mid-Generation đạt hiệu năng tổng thể cao nhất nhờ cấu trúc ba bước tường minh (trích xuất dữ kiện, phân tích, kết luận) giúp cải thiện độ trung thực. Pre-Generation thấp hơn một chút do không áp đặt cấu trúc suy luận, khiến mô hình dễ bỏ sót bước trung gian với câu hỏi phức tạp. Post-Generation cải thiện nhẹ Faithfulness nhờ lượt đối chiếu thứ hai, nhưng chi phí tính toán tăng gấp đôi do phải gọi mô hình hai lượt thay vì một. Kết quả xác nhận rằng một prompt được cấu trúc tốt ngay từ đầu mang lại hiệu quả cao hơn việc bổ sung bước kiểm tra sau cùng.

\subsubsection{So sánh với các hệ thống tham chiếu}

Bảng~\ref{tab:baseline_results} cho thấy GraphRAG vượt trội rõ rệt so với RAG truyền thống trên toàn bộ các độ đo IR: MAP 0,70 so với 0,50, NDCG@10 0,75 so với 0,55, Precision 0,80 so với 0,60 và Recall 0,65 so với 0,55 --- chênh lệch từ +0,10 đến +0,20 điểm. Việc khai thác cấu trúc đồ thị đưa các thực thể liên quan lên vị trí ưu tiên cao hơn trong kết quả truy xuất; trong khi RAG truyền thống trả về các đoạn văn bản có nội dung gần với câu hỏi nhưng rời rạc và thiếu liên kết logic, GraphRAG duy trì được mạch quan hệ giữa các thực thể thông qua cấu trúc đồ thị.

\begin{table}[h]
\centering
\caption{So sánh GraphRAG với các hệ thống tham chiếu trên bộ đánh giá CheoBench v2.}
\label{tab:baseline_results}
\begin{tabular}{lccc}
\hline
\textbf{Độ đo} & \textbf{GraphRAG} & \textbf{RAG truyền thống} & \textbf{LLM-only} \\
\hline
MAP            & 0.70 & 0.50 & -- \\
NDCG@10        & 0.75 & 0.55 & -- \\
Precision      & 0.80 & 0.60 & -- \\
Recall         & 0.65 & 0.55 & -- \\
\hline
\end{tabular}
\end{table}

LLM-only không có cơ chế truy xuất nên các độ đo truy xuất không áp dụng được; tuy nhiên ở khía cạnh tạo sinh, baseline này dễ gặp hiện tượng ảo giác khi gặp các thực thể và mối quan hệ đặc thù của miền Chèo --- xác nhận vai trò thiết yếu của cơ chế truy xuất có cấu trúc trong các miền tri thức chuyên biệt.

\subsection{Phân tích và thảo luận}

Kết quả thực nghiệm làm rõ hai nhóm phát hiện chính về kiến trúc GraphRAG đề xuất. Thứ nhất là vai trò then chốt của cấu trúc quan hệ đối với hiệu năng hệ thống; thứ hai là hiện tượng lợi ích biên giảm dần khi mở rộng ngữ cảnh. Phần này phân tích lần lượt từng nhóm dựa trên dẫn chứng định lượng và rút ra hàm ý thiết kế.

Vai trò của cấu trúc quan hệ thể hiện rõ ở hai khía cạnh so sánh. Cải thiện 0,22 điểm tổng hợp khi chuyển từ mức nút (0,54) sang kết hợp đa mức độ (0,76) cho thấy thông tin thuộc tính đơn lẻ không đủ cho các truy vấn có yếu tố quan hệ, còn chênh lệch +0,10 đến +0,20 so với RAG truyền thống trên toàn bộ các độ đo IR khẳng định ưu thế của truy xuất theo đồ thị so với truy xuất theo độ tương đồng văn bản. LLM-only càng bộc lộ hạn chế do không có nguồn tri thức để kiểm chứng, dễ gặp hallucination trong các câu hỏi đặc thù miền.

Hiện tượng lợi ích biên giảm dần quan sát được ở nhiều lát cắt: chênh lệch giữa tiểu đồ thị đơn lẻ (0,72) và kết hợp đa mức độ (0,76), cũng như giữa định dạng tốt nhất đơn lẻ (0,72) và kết hợp nhiều định dạng (0,76), đều chỉ khoảng 0,04 điểm; Post-Generation cũng chỉ cải thiện nhẹ so với Mid-Generation trong khi tốn gấp đôi chi phí gọi mô hình. Hiện tượng này phản ánh giới hạn cửa sổ chú ý của mô hình ngôn ngữ khi ngữ cảnh chứa nhiều thông tin không thiết yếu. Do đó, độ phong phú của ngữ cảnh cần được cân nhắc đồng thời với chi phí xử lý, thay vì cứ tăng kích thước là tốt hơn.

Tổng thể, kiến trúc GraphRAG phù hợp với miền tri thức có cấu trúc quan hệ phức tạp như nghệ thuật Chèo, giúp giảm phụ thuộc vào suy đoán nội tại của mô hình và tăng khả năng kiểm soát nguồn tri thức. Hạn chế chính còn lại nằm ở giới hạn cửa sổ ngữ cảnh và chi phí gọi mô hình ngôn ngữ. Các triển khai quy mô lớn hơn cần có cơ chế lựa chọn ngữ cảnh thông minh và tối ưu hoá chi phí suy luận.
